% =============================================================================
%  Notas de clase — Sesión 06: Práctica de Bit Manipulation
%  Programación Avanzada — Universidad Panamericana
%  Compilar: pdflatex sesion06_practica_bits.tex
% =============================================================================
\documentclass[12pt, oneside]{article}
\usepackage{progavanzada}

\begin{document}

\portada{Sesión 06}{Práctica de Bit Manipulation}{2026}

\tableofcontents
\newpage

% =============================================================================
\section{Técnicas fundamentales}
% =============================================================================

La sesión anterior cubrió los operadores de bits. Esta sesión los \textbf{aplicamos}.
Antes de pasar a los ejercicios, hay cuatro operaciones que aparecen en casi
todo código que manipula bits. Todas giran en torno a la misma idea:
\cpp{1 << k} es una \textbf{máscara de un bit} — solo el bit $k$ encendido.

\begin{center}
\begin{tabular}{llp{5.5cm}}
  \toprule
  \textbf{Operación} & \textbf{Expresión} & \textbf{Qué hace} \\
  \midrule
  Leer bit $k$    & \cpp{(x >> k) \& 1}
                  & Devuelve 1 si el bit $k$ está encendido, 0 si no \\[4pt]
  Encender bit $k$ & \cpp{x | (1 << k)}
                  & Pone el bit $k$ en 1 sin tocar los demás \\[4pt]
  Apagar bit $k$  & \texttt{x \& \textasciitilde(1 << k)}
                  & Pone el bit $k$ en 0 sin tocar los demás \\[4pt]
  Invertir bit $k$ & \cpp{x \^ (1 << k)}
                  & Cambia el bit $k$: 0 $\to$ 1, 1 $\to$ 0 \\
  \bottomrule
\end{tabular}
\end{center}

\begin{recuerda}
  \textbf{¿Por qué funcionan?}
  \begin{itemize}
    \item OR con \cpp{1 << k} \textbf{fuerza} el bit $k$ a 1 (OR con 1 siempre da 1).
    \item AND con \texttt{\textasciitilde(1 << k)} \textbf{fuerza} el bit $k$ a 0
          (AND con 0 siempre da 0).
    \item XOR con \cpp{1 << k} \textbf{invierte} solo el bit $k$
          (XOR con 1 cambia el valor; con 0 lo deja igual).
  \end{itemize}
\end{recuerda}

\begin{ejemplo}[\texttt{x = 0b10110100}, operaciones con distintos bits]
\begin{center}
\begin{tabular}{lll}
  \toprule
  \textbf{Operación} & \textbf{Resultado} & \textbf{Binario} \\
  \midrule
  \texttt{x} (original)           & 180 & \texttt{10110100} \\
  leer bit 1: \texttt{(x>>1)\&1}  & 0   & (el bit 1 está apagado) \\
  encender bit 1: \texttt{x|(1<<1)} & 182 & \texttt{10110110} \\
  apagar bit 2: \texttt{x\&\textasciitilde(1<<2)} & 176 & \texttt{10110000} \\
  invertir bit 7: \texttt{x\^{}(1<<7)} & 52 & \texttt{00110100} \\
  \bottomrule
\end{tabular}
\end{center}
\end{ejemplo}

% =============================================================================
\section{Máscaras de \textit{k} bits bajos}
% =============================================================================

Una operación frecuente es generar una máscara con los $k$ bits más bajos
encendidos y todos los demás apagados:

\begin{center}
\begin{tabular}{ccc}
  \toprule
  $k$ & \textbf{Máscara} & \textbf{Valor} \\
  \midrule
  3 & \texttt{00000111} &  7 \\
  5 & \texttt{00011111} & 31 \\
  8 & \texttt{11111111} & 255 \\
  \bottomrule
\end{tabular}
\end{center}

La expresión es \cpp{(1 << k) - 1}.

¿Por qué funciona? \cpp{1 << k} tiene exactamente el bit $k$ encendido
(vale $2^k$). Si le restamos 1, ese bit se apaga y todos los bits por
debajo de él se encienden — exactamente la máscara que necesitamos.

\begin{consejo}
  Combinar esta máscara con AND extrae los $k$ bits bajos de cualquier valor:
  \texttt{x \& ((1 << k) - 1)} equivale a \texttt{x \% (2\^{}k)},
  pero sin división.
\end{consejo}

% =============================================================================
\section{Ejercicios}
% =============================================================================

\subsection{Bloque A — Calentamiento}

Evalúa cada expresión \textbf{a mano}, mostrando los 8 bits de cada operando
en cada paso. Luego verifica con el cuaderno.

\begin{enumerate}

  \item \texttt{17 \& 12}
        \[
          17 = \texttt{????????}, \quad
          12 = \texttt{????????}, \quad
          \texttt{17 \& 12} = \texttt{????????} = \;?
        \]

  \item \texttt{7 << 2}
        \[
          7 = \texttt{????????}, \quad
          \texttt{7 << 2} = \texttt{????????} = \;?
        \]

  \item \texttt{13 | 12}
        \[
          13 = \texttt{????????}, \quad
          12 = \texttt{????????}, \quad
          \texttt{13 | 12} = \texttt{????????} = \;?
        \]

  \item \texttt{(\textasciitilde 7 | (1 << 1) | (1 << 3) | (1 << 5)) \& 29}

        \textit{Pista: calcula \texttt{\textasciitilde 7} primero (8 bits bajos),
        luego aplica cada OR, y al final el AND con 29.}

  \item \texttt{(\textasciitilde(13 >> 2)) \& 39}

        \textit{Este es el ejemplo resuelto en la sesión 05 — inténtalo sin mirar.}

\end{enumerate}

\subsection{Bloque B — Implementación}

\begin{enumerate}[resume]

  \item Escribe cuatro funciones en C++ que reciban \texttt{x} y \texttt{k},
        y apliquen cada técnica de la sección 1. Pruébalas con
        \texttt{x = 0b11001010} (202) y estos casos:

        \begin{center}
        \begin{tabular}{lll}
          \toprule
          \textbf{Operación} & \textbf{Caso} & \textbf{Esperado} \\
          \midrule
          Leer bit $k$      & $k=3$ & \texttt{1} \\
          Apagar bit $k$    & $k=6$ & \texttt{10001010} = 138 \\
          Encender bit $k$  & $k=0$ & \texttt{11001011} = 203 \\
          Invertir bit $k$  & $k=7$ & \texttt{01001010} = 74 \\
          \bottomrule
        \end{tabular}
        \end{center}

  \item Escribe una función \texttt{mascara(k)} que devuelva un entero con
        los $k$ bits más bajos encendidos, usando exactamente una expresión
        de bits. Verifica con $k \in \{3, 5, 8\}$.

\end{enumerate}

\subsection{Bloque C — Programación competitiva}

Los siguientes problemas tienen el formato de un juez en línea:
entrada, salida y restricciones. La solución debe ser eficiente:
\textbf{sin bucles donde baste una expresión de bits}.

\begin{enumerate}[resume]

  \item \textbf{¿Es potencia de 2?}

        \textbf{Entrada:} un entero $n$ con $1 \leq n \leq 10^9$. \\
        \textbf{Salida:} \texttt{1} si $n$ es potencia de 2, \texttt{0} si no.

        \begin{center}
        \begin{tabular}{cc}
          \textbf{Entrada} & \textbf{Salida} \\ \hline
          1  & 1 \\ 4 & 1 \\ 7 & 0 \\ 16 & 1 \\ 100 & 0
        \end{tabular}
        \end{center}

        \textit{Pista: una potencia de 2 tiene exactamente un bit encendido.
        ¿Qué hace \texttt{n \& (n-1)} con ese único bit?}

  \item \textbf{Cuenta bits encendidos (popcount).}

        \textbf{Entrada:} un entero $n$ con $0 \leq n \leq 10^9$. \\
        \textbf{Salida:} cantidad de bits en \texttt{1} en la representación
        binaria de $n$.

        \begin{center}
        \begin{tabular}{cc}
          \textbf{Entrada} & \textbf{Salida} \\ \hline
          0 & 0 \\ 7 & 3 \\ 13 & 3 \\ 255 & 8
        \end{tabular}
        \end{center}

        \textit{Pista A (bucle simple): usa \texttt{n \& 1} para leer el bit bajo
        y \texttt{n >>= 1} para avanzar.} \\
        \textit{Pista B (truco rápido): \texttt{n \& (n-1)} apaga el bit encendido
        de menor peso. Cuenta cuántas veces puedes repetirlo hasta llegar a 0.}

  \item \textbf{El elemento único.}

        \textbf{Entrada:} un arreglo donde todos los elementos aparecen
        \textbf{exactamente dos veces}, excepto uno. \\
        \textbf{Salida:} el elemento que aparece una sola vez. \\
        \textbf{Restricción:} tiempo $O(n)$, espacio $O(1)$ — sin arreglos auxiliares.

        \begin{center}
        \begin{tabular}{cc}
          \textbf{Arreglo} & \textbf{Salida} \\ \hline
          \{3, 5, 3, 7, 5\} & 7 \\
          \{1, 4, 1, 2, 4, 9, 2\} & 9
        \end{tabular}
        \end{center}

        \textit{Pista: recuerda que
        \texttt{x \textasciicircum{} x = 0} y
        \texttt{x \textasciicircum{} 0 = x}.
        ¿Qué pasa si aplicas XOR a todos los elementos del arreglo?}

  \item \textbf{Intercambio sin variable temporal.}

        Intercambia los valores de \texttt{a} y \texttt{b} usando
        \textbf{solo XOR}, sin declarar ninguna variable auxiliar.

        \textit{Pista: aplica XOR tres veces. Traza qué contiene cada variable
        después de cada paso.}

  \item \textbf{Distancia de Hamming.}

        La \textbf{distancia de Hamming} entre dos enteros es el número de
        posiciones de bit donde difieren.

        \textbf{Entrada:} dos enteros $a, b$ con $0 \leq a, b \leq 10^9$. \\
        \textbf{Salida:} su distancia de Hamming.

        \begin{center}
        \begin{tabular}{cc}
          \textbf{Entrada} & \textbf{Salida} \\ \hline
          $a=29,\ b=15$ & 2 \\
          $a=1,\ b=4$   & 2
        \end{tabular}
        \end{center}

        $29 = \texttt{11101}$, $15 = \texttt{01111}$,
        XOR $= \texttt{10010}$ — 2 bits distintos. \\[4pt]
        \textit{Pista: combina XOR con el popcount del problema anterior.}

  \item \textbf{El bit encendido más bajo.}

        \textbf{Entrada:} un entero $n > 0$. \\
        \textbf{Salida:} el \textbf{valor} del bit encendido de menor peso
        (el \texttt{1} más a la derecha).

        \begin{center}
        \begin{tabular}{cc}
          \textbf{Entrada} & \textbf{Salida} \\ \hline
          12 (\texttt{1100}) & 4 \\
          24 (\texttt{11000}) & 8 \\
          7  (\texttt{0111}) & 1 \\
          40 (\texttt{101000}) & 8
        \end{tabular}
        \end{center}

        \textit{Pista: en complemento a 2, \texttt{-n} invierte todos los bits y
        suma 1. ¿Qué produce \texttt{n \& -n}?}

  \item \textbf{Paridad.}

        \textbf{Entrada:} un entero $n \geq 0$. \\
        \textbf{Salida:} \texttt{0} si la cantidad de bits en \texttt{1}
        es par, \texttt{1} si es impar.

        \begin{center}
        \begin{tabular}{cc}
          \textbf{Entrada} & \textbf{Salida} \\ \hline
          7 (\texttt{111}) & 1 \\
          6 (\texttt{110}) & 0 \\
          255 (\texttt{11111111}) & 0 \\
          13 (\texttt{1101}) & 1
        \end{tabular}
        \end{center}

        \textit{Pista: la paridad es el XOR de todos los bits del número.
        Puedes acumular con \texttt{res \textasciicircum{}= (n \& 1)}
        y desplazar.}

  \item \textbf{★ Contar subconjuntos (bitmask).}

        Una \textbf{máscara de bits} de $k$ bits representa todos los
        subconjuntos posibles de un conjunto de $k$ elementos: cada bit
        $i$ indica si el elemento $i$ está incluido.

        \textbf{Tarea:} dado $k$, imprime todas las máscaras de 0 a
        $2^k - 1$, mostrando qué elementos incluye cada una.

        Para $k = 3$ (elementos \{A, B, C\}):
        \begin{verbatim}
000  ->  {}
001  ->  {A}
010  ->  {B}
011  ->  {A, B}
100  ->  {C}
101  ->  {A, C}
110  ->  {B, C}
111  ->  {A, B, C}
        \end{verbatim}

        \textit{Pista: itera \texttt{mask} de 0 a \texttt{(1<<k)-1}.
        Para cada \texttt{mask}, verifica cada bit con
        \texttt{(mask >> i) \& 1}.}

\end{enumerate}

% =============================================================================
\section{Resumen de técnicas}
% =============================================================================

\begin{recuerda}
\begin{center}
\begin{tabular}{lll}
  \toprule
  \textbf{Técnica} & \textbf{Expresión} & \textbf{Uso en competencia} \\
  \midrule
  Leer bit $k$       & \cpp{(x >> k) \& 1}           & Verificar si un elemento está en subconjunto \\
  Encender bit $k$   & \cpp{x | (1 << k)}             & Agregar elemento a subconjunto \\
  Apagar bit $k$     & \texttt{x \& \textasciitilde(1 << k)} & Quitar elemento de subconjunto \\
  Invertir bit $k$   & \cpp{x \^ (1 << k)}            & Toggle de estado \\
  Máscara $k$ bits   & \cpp{(1 << k) - 1}             & Extraer los $k$ bits bajos \\
  ¿Potencia de 2?    & \cpp{n > 0 \&\& !(n \& (n-1))} & Verificar en $O(1)$ \\
  Bit más bajo       & \cpp{n \& -n}                  & Iterar sobre bits encendidos \\
  Apagar bit más bajo & \cpp{n \& (n-1)}              & Popcount eficiente \\
  Elemento único     & XOR de todo el arreglo         & Problema clásico $O(n)$/$O(1)$ \\
  Distancia de Hamming & \cpp{popcount(a \^ b)}        & Códigos de error, similitud \\
  Iterar subconjuntos & \cpp{for mask in 0..(1<<k)-1}  & DP sobre subconjuntos \\
  \bottomrule
\end{tabular}
\end{center}
\end{recuerda}

% =============================================================================
\section{Metacognición}
% =============================================================================

Escribe un \textbf{correo electrónico} a tu papá, mamá o tutor contándole
sobre los temas que hemos visto en clase hasta el momento.

El correo debe incluir:

\begin{itemize}
  \item \textbf{Qué sabías} del tema antes de esta clase.
  \item \textbf{Qué no sabías} y aprendiste.
  \item \textbf{Qué te gustó} de la actividad.
  \item \textbf{Cómo piensas aplicarlo} en tu vida profesional.
\end{itemize}

Escríbelo como si le estuvieras \textit{contando} a esa persona, con tus
propias palabras, no como un reporte técnico.  Pon al profesor en copia
(\textbf{cc}).

\end{document}
